\documentclass[../../../include/open-logic-section]{subfiles}

\begin{document}

\olfileid[fa-IR]{sfr}{infinite}{card-sb}

\olsection{ضمیمه: اثبات شرودر–برنشتاین}

پیش از آنکه نظریهٔ طبیعی مجموعه‌ها را ترک کنیم، باید یک برهان طبیعیِ نهایی
(اما ظریف!) را بررسی کنیم. این، برهانِ قضیهٔ شرودر–برنشتاین
(\olref[sfr][siz][sb]{thm:schroder-bernstein}) است: اگر
$\cardle{A}{B}$ و $\cardle{B}{A}$، آنگاه $\cardeq{A}{B}$؛ یعنی، با داشتن
!!{injection}s $f \colon A \to B$ و $g \colon B \to A$، !!a{bijection}
$h \colon A \to B$ وجود دارد.

در این فصل، از مفهوم \emph{بستارها} نزد ددکیند پیروی کردیم. در واقع،
ددکیند با استفاده از این مفهوم برهانی دل‌انگیز برای شرودر–برنشتاین عرضه
کرد و ما آن را اینجا ارائه خواهیم کرد. برهان حاضر از نزدیک از
\citet[صص.~157--8]{Potter2004} پیروی می‌کند؛ اگر شرحی اندکی متفاوت اما در
اساس مشابه می‌خواهید، به آن بنگرید. اندکی جست‌وجو در وب نیز شما را متقاعد خواهد کرد که این
قضیه---همچون گنگ‌بودنِ $\sqrt{2}$---از آن قضیه‌هایی است که برهان‌های جالب
و متفاوتِ \emph{بسیاری} برایشان وجود دارد.

با استفاده از نمادگذاری‌ای مشابه آنچه در \olref[sfr][infinite][dedekind]{Closure} آمده است،
قرار دهید
\[
\Closureofunder{f}{B} = \bigcap \Setabs{X}{B \subseteq X
\text{ و $X$ تحت $f$ بسته است}}
\]
برای هر مجموعهٔ $B$ و تابع~$f$. با این تعریف، $\Closureofunder{f}{B}$
کوچک‌ترین مجموعهٔ بسته تحت $f$ است
که~$B$ را دربرمی‌گیرد، بدین معنا که:

\begin{lem}\ollabel{Closureprops}
	برای هر تابع $f$ و هر $B$:
	\begin{enumerate}
		\item\ollabel{Closurehaselem} $B \subseteq \Closureofunder{f}{B}$؛ و
		\item\ollabel{Closureclosed} $\Closureofunder{f}{B}$ تحت $f$ بسته است؛ و
		\item\ollabel{Closuresmallest} اگر $X$ تحت $f$ بسته باشد و $B
		\subseteq X$، آنگاه $\Closureofunder{f}{B} \subseteq X$.
	\end{enumerate}
\end{lem}

\begin{proof}
دقیقاً همانند \olref[sfr][infinite][dedekind]{closureproperties}.
\end{proof}

برای رسیدن به شرودر–برنشتاین به یک واقعیت دیگر نیاز داریم:

\begin{prop}\ollabel{sbhelper}
اگر $A \subseteq B \subseteq C$ و $A \approx C$، آنگاه $\cardeq{\cardeq{A}{B}}{C}$.
\end{prop}

\begin{proof}
با داشتن !!a{bijection} $f \colon C \to A$، بگذارید $F =
\Closureofunder{f}{C \setminus B}$، و تابعی چون $g$ با دامنهٔ $C$ به‌صورت
زیر تعریف کنید:
\[
	g(x) =
	\begin{cases}
			f(x) &\text{اگر $x \in F$}\\
			x & \text{در غیر این صورت}
		\end{cases}
\]
نشان خواهیم داد که $g$ یک !!a{bijection} به‌صورت $C \to B$ است؛ از این امر
نتیجه خواهد شد که $\comp{f^{-1}}{g} \colon A \to B$ یک !!a{bijection}
است، و برهان کامل خواهد شد.

نخست ادعا می‌کنیم اگر $x \in F$ اما $y\notin F$، آنگاه $g(x) \neq
g(y)$. برای برهان خلف، خلاف آن را فرض کنید؛ در این صورت $y = g(y) = g(x) =
f(x)$. چون $x \in F$ و بنا بر \olref{Closureprops}، $F$ تحت $f$ بسته
است، داریم $y = f(x) \in F$، که تناقض است.

اکنون فرض کنید $g(x) = g(y)$. پس، بنا بر آنچه در بالا آمد، $x \in F$ اگر
و تنها اگر $y \in F$. اگر $x, y \in F$، آنگاه $f(x) = g (x) = g(y) = f(y)$؛
پس $x = y$، زیرا $f$ یک !!a{bijection} است. اگر $x, y \notin F$، آنگاه
$x = g(x) = g(y) = y$. بنابراین $g$ یک !!a{injection} است.

تنها مانده است نشان دهیم $\ran{g} = B$. پس $x \in B \subseteq C$ را ثابت
بگیرید. اگر $x \notin F$، آنگاه $g(x) = x$. اگر $x \in F$، آنگاه
$x = f(y)$ برای یک $y \in F$؛ زیرا در غیر این صورت $F \setminus \{x\}$
تحت $f$ بسته می‌بود و $C\setminus B$ را دربرمی‌گرفت، که بنا بر
\olref{Closureprops} ناممکن است؛ اکنون $g(y) = f(y) = x$.
\end{proof}

سرانجام، این هم برهان نتیجهٔ اصلی. به یاد آورید که برای تابع $h$ و مجموعهٔ
$D$، تعریف می‌کنیم $\funimage{h}{D} = \Setabs{h(x)}{x
\in D}$.

\begin{proof}[برهان شرودر–برنشتاین] فرض کنید $f \colon A \to B$
و $g \colon B \to A$، !!{injection}s باشند. چون $\funimage{f}{A}
\subseteq B$، داریم $\funimage{g}{\funimage{f}{A}} \subseteq
g[B] \subseteq A$. همچنین، $\comp{f}{g} \colon A \to
\funimage{g}{\funimage{f}{A}}$ یک !!{injection} است، زیرا $g$ و
$f$ هر دو چنین‌اند؛ و در واقع $\comp{f}{g}$ یک !!a{bijection} است، صرفاً
به سبب نحوه‌ای که هم‌دامنهٔ آن را تعریف کردیم. پس
$\cardeq{\funimage{g}{\funimage{f}{A}}}{A}$، و ازاین‌رو بنا بر
\olref{sbhelper}، یک !!a{bijection} چون $h \colon A \to
\funimage{g}{B}$ وجود دارد. افزون بر این، $g^{-1}$ یک !!a{bijection}
$\funimage{g}{B} \to B$ است. بنابراین $\comp{h}{g^{-1}} \colon A \to B$
یک !!a{bijection} است.
\end{proof}

\end{document}
